% ============================================================
%  给韩尧宇学弟的一封信 —— 正式中文书信
%  编译方式: XeLaTeX
%  特色: 信头 + 装饰双线 + 精致落款 + 正文分点
% ============================================================
\documentclass[12pt,a4paper]{article}

% ---------- 页面设置 ----------
\usepackage[margin=2.8cm, top=3cm, bottom=3cm]{geometry}
\usepackage{ctex}
\usepackage{xcolor}
\usepackage{setspace}
\usepackage{tcolorbox}
\usepackage{enumitem}
\usepackage{fancyhdr}

% ---------- 行距 ----------
\onehalfspacing

% ---------- 页眉：极简信纸感 ----------
\setlength{\headheight}{14pt}
\pagestyle{fancy}
\fancyhf{}
\renewcommand{\headrulewidth}{0.4pt}
\fancyhead[C]{\small\color{gray}缄 \  ·  \ 【2026年8月31日】}
\fancyfoot[C]{\small\color{gray}—— 第 \thepage 页 ——}

% ---------- 颜色 ----------
\definecolor{ink}{RGB}{60,60,60}
\definecolor{deepblue}{RGB}{30,60,110}
\definecolor{softgold}{RGB}{180,150,80}
\definecolor{lightbg}{RGB}{246,243,238}

\begin{document}
\color{ink}

% ============================================================
%	信头
% ============================================================
\begin{tcolorbox}[colback=lightbg, colframe=deepblue, boxrule=1pt,
                  arc=2mm, left=6mm, right=6mm, top=4mm, bottom=4mm]
  \centering
  {\color{deepblue}\Large\heiti 别只顾着赶路，也看看路边的风景}\\[0.4em]
  {\color{deepblue}\normalsize ——给即将面试志道班的你}\\[0.6em]
  \textcolor{softgold}{\rule{0.4\textwidth}{1pt}\ $\diamond$\ \rule{0.4\textwidth}{1pt}}\\[0.5em]
  {\small\color{gray}【2026年8月31日】}
\end{tcolorbox}

\vspace{1.4em}

% ============================================================
%	称呼
% ============================================================
\noindent\textbf{\large 韩尧宇学弟如晤}：\par
\vspace{0.4em}

% ============================================================
%	正文
% ============================================================
深夜提笔，窗外是南京的秋。紫金山的轮廓在月光下像一匹安静伏卧的兽，风穿过梧桐叶，沙沙的，像谁在翻一本很旧的书。\par

而你在江阴。申港河边。\par

我查过那条河。北接长江，南入西横河，是春申君当年开凿的申浦故道。两千多年了，水还在流，不急不缓，拐过申港镇的街巷，拐过你教学楼外的那排香樟，拐过你每天匆匆走过的那座桥。\par

你有没有在傍晚停下来看过它？\par

不是那种"特意去观察"的看，而是某天从自习室出来，天已经暗了，路灯一盏盏亮起来，河面上浮着碎金一样的水光。那一刻你站在桥头，什么都没想，只是看着——然后心里忽然软了一下。\par

我想说的就是这个"软"。\par

我觉得老师要找的，不是最硬的剑，是心里还留着这块软地方的人。\par

咱俩都是海安人。里下河平原上长大的孩子，从小看惯了水。串场河不宽，水流不急，河岸边长着不知名的草，夏天有蜻蜓停在上面，翅膀是透明的。那种地方教给你的东西，和任何一本教材都不一样：它教你慢，教你等，教你相信水终究会流到它该去的地方。\par

可你到了大学，到了江阴校区，到了要面试志道班的时候，我觉得你可能有点急了。\par

我理解，我当年也这样。周围的人都在跑，我不敢停。你开始背钱学森的故事、背"系统工程""拔尖创新""国家战略"这些词。你把它们排练了一遍又一遍，像在准备一场演出。\par

但学弟，我想问你一个很私人的问题——\par

你上一次被什么东西真正打动，是什么时候？\par

不是被"优秀学长分享会"打动，不是被"保研率"打动，而是被一个具体的、微小的、说不清道不明的瞬间打动。可能是某天傍晚申港天边的晚霞，可能是某个实验课上仪器突然亮起来的那一下，可能是深夜耳机里某句歌词恰好撞上了你的心事。\par

那个瞬间，才是你的"道"。\par

老师不傻。他们见过太多把"志存高远"挂在嘴边的人，也见过太多眼睛里没有光的人。他们想看到的，是你提到某件事的时候，声音会不自觉地低下来，像在说一个秘密。他们想找的，是一个具有一点天赋的"人"。\par

所以别演。这是唯一的，也是不算技巧的技巧。我当年的技巧，我问过了我们班上所有的同学，他们也是如此。\par

别演。\par

我可以和你交流一下我当时面试的情况，我面试的是鼎新班的芮筱亭院士团队和李鸿志院士团队。\par

第一个面试的是李鸿志院士团队。当时他们问了我两个问题：第一个比较专业，是问我关于主动防御和被动防御的看法和理解。我当时想到啥说啥，也记不清了。第二个问题是问我为什么把他们放在第二志愿。我答的是我只想了解兵器，无论兵器的哪个方面，先后没怎么在意。\par

第二个面试是芮筱亭院士团队，是团队的王燕老师问的问题。第一个问题是让我说一下是哪个高中毕业的，江苏省海安高级中学是加分项。第二个问题是问我平时有什么爱好，我按实际情况说了。不知道你后面会不会和我有交集，我这人兴趣爱好挺广泛的，我现在是学院话剧团的导演，学弟有兴趣可以和我交流这个。\par

这些是我知道的一些信息。\par

过几天你要坐大巴来南京面试。过江的时候，别低头看资料。抬头看看窗外——长江在这里收得很窄，水被挤得急了些，浪花翻着白。可过了这一段，江面会忽然铺开，宽得看不见对岸。\par

人也是这样。有些窄，是为了让你学会在逼仄里站稳。有些宽，是留给你的。\par

我在南京等你。不管结果怎样，来孝陵卫找我。我带你去食堂吃鸭血粉丝（资金链有点紧张），带你去看日落，带你在三号路上走走。你会明白——有些路，不是选出来的，是走着走着，自己就清楚了。\par

申港河不急。你也不急。\par

祝好。

% ============================================================
%	落款
% ============================================================
\vspace{1.6em}
\begin{flushright}
  {\heiti 魏舟}\\[0.5em]
\end{flushright}

\vspace{2.5em}
\begin{center}
  \textcolor{softgold}{\rule{0.6\textwidth}{0.6pt}}\\[0.3em]
  {\footnotesize\color{gray} — 江声未老，来日方长 —}
\end{center}

\end{document}
